
\section*{3. Phenomena specific to characteristic \(0\)}
\phantomsection
\label{sga3:VIIB:3}
\addcontentsline{toc}{section}{3. Phenomena specific to characteristic 0}

Throughout \SGARefLink{sga3:VIIB:3}{Section~3}, we assume that the pseudocompact ring \(k\)
contains the field of rational numbers \(\mathbb Q\).

\subsection*{3.1. Lemma}
\phantomsection
\label{sga3:VIIB:3.1}

\noindent\textbf{Lemma 3.1.}
\phantomsection
\label{sga3:VIIB:3.1:lemma}
{\itshape Let \(C\) be a commutative unital
\(\mathbb Q\)-algebra, and let \(L\) be a Lie algebra over \(C\)
whose underlying \(C\)-module is free. Then the canonical map
\[
  L\longrightarrow U(L)
\]
identifies \(L\) with the set of primitive elements of \(U(L)\).}

\textit{Proof.}
Identify \(L\) with its canonical image in \(U(L)\). Let \(I\) be
a totally ordered set, and let
\[
  (x_i)_{i\in I}
\]
be a basis of \(L\) indexed by \(I\). Denote by
\(\mathbb N^{(I)}\) the set of families
\[
  n=(n_i)_{i\in I}
\]
of nonnegative integers for which \(n_i=0\) except possibly at
finitely many indices
\[
  i_1<i_2<\cdots<i_s;
\]
these indices depend on \(n\). Put
\[
  x^n
  =
  x_{i_1}^{\,n_{i_1}}
  x_{i_2}^{\,n_{i_2}}
  \cdots
  x_{i_s}^{\,n_{i_s}},
  \qquad
  n!
  =
  (n_{i_1}!)(n_{i_2}!)\cdots(n_{i_s}!).
\]
By the Poincaré--Birkhoff--Witt theorem, the elements \(x^n\)
form a basis of \(U(L)\). Moreover,
\begin{equation}
  \Delta_{U(L)}
  \left(\frac{x^n}{n!}\right)
  =
  \sum_{0\leq m\leq n}
  \frac{x^m}{m!}
  \otimes
  \frac{x^{\,n-m}}{(n-m)!}.
  \tag{*}
  \label{sga3:VIIB:3.1:pbw-comultiplication}
\end{equation}
Here the sum ranges over all \(m\in\mathbb N^{(I)}\) such that
\(0\leq m\leq n\), that is,
\[
  0\leq m_i\leq n_i
  \qquad\text{for every }i.
\]
It follows immediately that
\[
  \Delta_{U(L)}(u)=u\otimes1+1\otimes u
\]
if and only if \(u\) is a linear combination of the \(x_i\).
This proves the lemma.
